\begin{abstract}
Markets for data promise to unlock its economic value, yet in practice they remain far less active than expected---largely because those who supply data are not rewarded for the value it creates downstream. A defining feature of data is its \emph{replicability}: a buyer can refine purchased data into a new product and resell it to \emph{many} downstream buyers, so a single source seeds a branching cascade of resales that naturally forms a \emph{tree}. Because an upstream seller captures none of this downstream value, its incentive to trade is weakened. Existing works propose \emph{profit reallocation}---returning part of downstream revenue to upstream contributors---as a natural remedy. But whether profit reallocation works on the tree-structured markets that replicable data actually induces has remained open.

To bridge this gap, we develop a principled framework for profit reallocation on tree-structured data markets. We introduce a sequential trading game on a tree and a general class of budget-feasible profit reallocation mechanisms (PRMs) over it. We derive efficient algorithms to compute the induced equilibria---a polynomial-time exact algorithm for discrete valuations and a fully polynomial-time approximation scheme (FPTAS) for continuous ones---via a subtree decomposition technique that tames the potential coupling across a seller's children. We then prove that, under mild assumptions, \emph{any} budget-feasible PRM weakly expands the trades that occur in equilibrium, and any budget-balanced PRM additionally weakly improves social welfare, relative to the baseline that reallocates nothing. Experiments on synthetic markets confirm that these benefits are substantial, persist even when the assumptions fail, and grow with the depth and branching of the tree.

\keywords{Data Trading \and Profit Reallocation Mechanism \and Equilibrium Analysis \and Tree Structure}
\end{abstract}
